\documentclass{situationist-preprint}

\title{Moral Twin Earth Without Twin Meanings}
\msid{MTE-13-a}
\msfamily{MTE-13}
\msdiscipline{Metaethics}
\mstype{Principal manuscript}
\mskeywords{Moral Twin Earth; ethical naturalism; reference; causal theory of reference; moral disagreement; meaning change}

\begin{document}
\maketitle

\begin{abstract}
Horgan and Timmons's Moral Twin Earth describes a community whose use of \enquote{wrong} is causally regulated by a deontological property while ours is regulated by a consequentialist one. On the causal-regulation semantics of Cornell realism, the two words refer to different properties; yet if the communities disputed whether an act is wrong, they would, intuitively, genuinely disagree. Horgan and Timmons conclude that no naturalist semantics can capture the meaning of moral terms. We argue that moral terms refer to the same property on Twin Earth, and that the divergence in use reflects error, not a difference in meaning. The reference of a moral term is fixed by the role that moral inquiry assigns to its referent, not by whatever causally regulates the term's use. The Twin Earthers' inquiry assigns their term the same role as ours, and the role has a unique realizer, the convergence of idealized reactions, which is a natural property. We show that this semantics predicts the disagreement intuition, gives a principled criterion separating error from meaning change, and is naturalist, so that Horgan and Timmons's generalization to all naturalist semantics fails.
\end{abstract}

\section{The problem}

Putnam's Twin Earth is a planet like Earth except that the liquid in its lakes and rivers, which its inhabitants call \enquote{water}, is not H$_2$O but a different substance, XYZ, which is superficially indistinguishable from it \citep{Putnam1975}. Most who consider the case judge that the Twin Earthers' word does not mean what ours means, since it refers to XYZ. If an Earthling and a Twin Earthling disputed whether a sample of H$_2$O is \enquote{water}, they would be talking past each other. The case supports a semantics on which the reference of natural kind terms is fixed by the underlying nature of the samples to which their users are causally related, and LNG-08 draws from it the conclusion that meanings are not in the head, though they depend on what is in the head.

Naturalist moral realists of the Cornell school applied a semantics of this kind to moral terms. On the view of \citet{Boyd1988}, a moral term refers to the natural property that causally regulates its use: the property such that the community's tendency to apply the term is, over time and under the influence of evidence, responsive to the presence of that property. The view made moral terms resemble natural kind terms, and allowed that moral properties are natural properties though no moral term is synonymous with any naturalistic description.

\citet{HorganTimmons1991} described a Moral Twin Earth. Its inhabitants use \enquote{good} and \enquote{wrong} as we do in every functional respect: the words guide their actions, express their approval, figure in their deliberation and are subject to argument and revision. But their use of the words is causally regulated by a deontological property, while ours, we may suppose, is regulated by a consequentialist one. On the causal-regulation semantics, their words refer to the deontological property and ours to the consequentialist one, so the words differ in meaning. Yet, Horgan and Timmons observe, if a group of Earthlings and Twin Earthlings disputed whether some act is wrong, it would seem that they genuinely disagree, not that they talk past each other. The case of water and the case of morality elicit opposite intuitions. Horgan and Timmons concluded that the causal-regulation semantics is false, and argued in later work that the argument generalizes to any semantics on which moral terms refer to natural properties in virtue of a natural relation between their use and the properties \citep{HorganTimmons1992a,HorganTimmons1992b}.

Responses have taken several forms. \citet{Copp2000} argued that the case is underdescribed, and that once the details are filled in it is unclear that the intuition of disagreement survives. \citet{Merli2002} proposed that moral terms refer to the properties that would regulate their use under idealized conditions, so that the actual regulation of use is not decisive. \citet{vanRoojen2006} argued that the intuition of disagreement reflects our knowledge that the two communities share enough about the role of moral terms to disagree. \citet{PlunkettSundell2013} argued that the communities may disagree even if their terms differ in meaning, by engaging in a metalinguistic negotiation about how the terms should be used. \citet{Dowell2016} questioned whether the intuitions of ordinary speakers about such cases have the evidential significance that the argument requires.

We argue for the conclusion that the intuition of disagreement is correct and that the semantics of moral terms vindicates it. The reference of moral terms is fixed by role. This is the semantics of MTE-01, on which moral facts are natural facts about idealized reactions, and its application to Moral Twin Earth is the subject of this manuscript.

\section{Framework}

\begin{definition}[Causal-regulation semantics]\label{d:crs}
On the \emph{causal-regulation semantics}, a moral term $t$ of a community refers to the natural property that causally regulates the community's use of $t$.
\end{definition}

\begin{definition}[Role semantics]\label{d:rs}
On the \emph{role semantics}, a moral term $t$ of a community refers to the property that realizes the role the community's inquiry assigns to the referent of $t$. For \enquote{wrong}, the role has five clauses: its referent supervenes on the natural features of acts; a rational agent who believes that an act has it is, insofar as she is rational, disposed not to perform it; moral judgment under improving epistemic conditions tends toward the truth about it; it depends on how an act bears on everyone affected; and any actual agent or community can be mistaken about it.
\end{definition}

The role semantics is Principle~P1 of MTE-01, with the role given by Definition~3.4 there. MTE-01 argues that the role has a unique realizer up to coextension (Assumption~3.6 there), and that it is realized by idealized disapproval, the convergence of the idealized reactions of all agents on disapproval. The role semantics is a semantics of the kind LNG-02 defends for names, on which reference is causal with a descriptive component that takes over when the causal chain fails to determine a referent; for moral terms, the descriptive component is the role, and, as the present case shows, the causal component does not determine the referent.

\begin{definition}[Moral Twin Earth]\label{d:mte}
A \emph{Moral Twin Earth} is a community such that:
\begin{enumerate}[label=(\alph*)]
\item its use of \enquote{wrong} is functionally like ours: the word guides action, is used to express disapproval, figures in deliberation, is applied impartially, and is subject to argument, evidence and revision;
\item its use of \enquote{wrong} is causally regulated by a property $G_T$ distinct from the property $G_E$ that regulates ours;
\item its members are like us in nature: in their capacities for welfare and suffering, for sympathy and for reasoning.
\end{enumerate}
\end{definition}

Clauses (a) and (b) are Horgan and Timmons's stipulations. Clause (c) is implicit in their description, on which the Twin Earthers are as much like us as the scenario permits; we make it explicit because it matters for the analysis, and we consider in Section~\ref{s:cases} the variant in which it fails.

\begin{principle}[Disagreement datum]\label{p:datum}
If an Earthling asserts and a Twin Earthling denies that a given act is wrong, they genuinely disagree.
\end{principle}

We accept Principle~\ref{p:datum}. Horgan and Timmons are right that the intuition is robust, and we take it to be evidence about the semantics of moral terms.

\section{The argument}

\begin{lemma}\label{l:role}
The Twin Earthers' inquiry assigns to the referent of their \enquote{wrong} the same role as our inquiry assigns to the referent of ours.
\end{lemma}

\begin{proof}
The five clauses of the role in Definition~\ref{d:rs} are features of the functional role of the term in a community's practice: whether its application is constrained by natural features, whether it guides the action of those who accept it, whether it is subject to correction by evidence and argument, whether it is applied impartially, and whether its users treat themselves as fallible about it. By clause (a) of Definition~\ref{d:mte}, the Twin Earthers' use of \enquote{wrong} is functionally like ours in each of these respects. Nothing in the role concerns which property regulates the use; that is the respect in which, by clause (b), the communities differ.
\end{proof}

\begin{theorem}\label{t:same}
On the role semantics, \enquote{wrong} on Earth and \enquote{wrong} on Moral Twin Earth refer to the same property.
\end{theorem}

\begin{proof}
By Lemma~\ref{l:role}, the two terms are assigned the same role. By the role semantics, each refers to the realizer of that role. The role has a unique realizer up to coextension, which by MTE-01 is idealized disapproval: the property an act has when the idealized reactions of all agents to it converge on disapproval. The property is not relative to a community, since it concerns the idealized reactions of all agents, and by clause (c) of Definition~\ref{d:mte} the Twin Earthers' idealized reactions do not differ from ours. So both terms refer to idealized disapproval.
\end{proof}

\begin{corollary}\label{c:error}
If an Earthling asserts and a Twin Earthling denies that a given act is wrong, at least one of them is mistaken.
\end{corollary}

\begin{proof}
By Theorem~\ref{t:same}, the assertion and the denial concern the same property of the same act, and so cannot both be true.
\end{proof}

\begin{proposition}\label{p:explains}
The role semantics predicts Principle~\ref{p:datum}; the causal-regulation semantics contradicts it.
\end{proposition}

\begin{proof}
Two speakers genuinely disagree when one asserts and the other denies the same proposition. On the role semantics, by Theorem~\ref{t:same}, the Earthling's assertion and the Twin Earthling's denial concern the same proposition. On the causal-regulation semantics, the Earthling asserts that the act has $G_E$ and the Twin Earthling denies that it has $G_T$, which are different propositions, and both may be true.
\end{proof}

Theorem~\ref{t:same} and Corollary~\ref{c:error} establish the conclusion of this manuscript: moral terms refer to the same properties on Twin Earth, and the apparent divergence reflects error. The divergence in what regulates use is a divergence in which features of acts the two communities respond to when they apply \enquote{wrong}. At most one of these responses tracks idealized disapproval, and it may be that neither does: perhaps the consequentialist property and the deontological property each coincide with idealized disapproval over part of their extension and diverge from it over the rest, in which case both communities are partly in error.

\section{Error and meaning change}\label{s:cases}

A semantics for moral terms must distinguish disagreement from talking past one another, and therefore error from difference in meaning. The role semantics gives a criterion.

\begin{proposition}[Criterion]\label{p:crit}
Two communities' terms differ in meaning if and only if their inquiries assign different roles to the referents of the terms. Where the roles are the same, a divergence in application is error on at least one side.
\end{proposition}

\begin{proof}
On the role semantics, reference is a function of role. If the roles are the same, the referents are the same, and divergent applications cannot all be correct. If the roles differ, the referents may differ, since a different role may have a different realizer.
\end{proof}

We apply the criterion to four variants.

\paragraph{The original case.} Horgan and Timmons's Twin Earthers assign \enquote{wrong} the same role as we do, by clause (a) of Definition~\ref{d:mte}, so their term means what ours means, and the divergence is error. This agrees with the intuition.

\paragraph{The etiquette planet.} Imagine a community whose word \enquote{wrong} is applied to breaches of local custom, is not subject to correction by argument, is not taken to concern everyone affected, and whose users do not regard their community as possibly mistaken about it. The role of this term lacks clauses (c), (d) and (e). By Proposition~\ref{p:crit}, it may differ in meaning from our \enquote{wrong}, and an Earthling and a member of this community who disputed whether an act is \enquote{wrong} would be talking past each other. This, too, agrees with the intuition, which in this case is that the community's word is closer to our \enquote{improper} or \enquote{not done}. The role semantics thus does not make every evaluative word of every community mean what our \enquote{wrong} means; it makes it do so exactly when the word plays the role of a moral term.

\paragraph{Twin Earthers of another nature.} Suppose that clause (c) of Definition~\ref{d:mte} fails: the Twin Earthers cannot suffer, or have capacities for welfare very different from ours. Their term \enquote{wrong} still plays the same role, and still refers to idealized disapproval. But the idealized reactions to acts that affect them differ from the idealized reactions to the corresponding acts that affect us, because the acts affect beings of a different nature. The extension of the one property then varies with the nature of those concerned. This is the conclusion of MTE-02: moral truths hold for all rational agents, but their content is relative to the nature of the agents concerned. The Twin Earthers who apply \enquote{wrong} differently to acts affecting themselves need not be in error, since the acts differ in their natural features; and they mean what we mean.

\paragraph{Divergence under idealization.} Suppose that, on some question, the idealized reactions of Earthlings and Twin Earthlings would themselves diverge. Then, on MTE-01, the act lacks idealized disapproval, and so is not wrong; its omission is not wrong either, and the comparative question which option is better has no answer. Since the two communities' terms refer to the same property, they still disagree, and an Earthling who asserts that the act is wrong is mistaken, as is a Twin Earthling who asserts that it is required. Divergence under idealization yields permissions and incomparabilities, not a difference in meaning.

The four cases show that the role semantics draws the line between error and meaning change where our judgments draw it, and that it explains why Twin Earth for moral terms elicits the opposite intuition to Twin Earth for natural kind terms.

\section{Objections and replies}

\begin{objection}[Generalization]
Horgan and Timmons argue that the Moral Twin Earth argument generalizes: any naturalist semantics fixes the reference of moral terms by some natural relation between their use and a property, and for any such relation one can construct a Twin Earth on which the relation picks out a different property while the intuition of disagreement persists.
\end{objection}

\begin{reply}
The role semantics is naturalist: the role is a feature of a community's practice, its realizer is a natural property, and the relation of realizing a role is a natural relation. Yet no Twin Earth of the kind the objection requires can be constructed for it. To make the role semantics assign the Twin Earthers' term a different referent, one must give their term a different role. But a community whose moral term plays a different role, lacking the inquiry, impartiality or fallibility clauses, is the etiquette planet, and there the intuition of disagreement does not persist. The generalization assumes that every naturalist relation of reference is a relation between use and the property that regulates use. The role semantics shows that the assumption is false.
\end{reply}

\begin{objection}[The water analogy]
The role semantics would also give the wrong result for \enquote{water}. Twin Earthers' \enquote{water} plays the same role in their practice as ours: it names the clear, drinkable liquid of lakes and rivers. On the role semantics, it would refer to the same thing as ours, but it refers to XYZ.
\end{objection}

\begin{reply}
The role of a natural kind term includes a clause that moral terms lack: deference to the underlying nature of local paradigm samples. Users of \enquote{water} take the word to apply to whatever has the same nature as the stuff in their lakes, and it is this clause that makes the reference of \enquote{water} vary with environment, as \citet{Putnam1975} argued. Users of \enquote{wrong} take no such attitude toward paradigm cases. By clause (e) of the role, they regard any community, their own included, as possibly mistaken even about the acts it treats as paradigms; slavery was a paradigm of permissible conduct for many communities that were mistaken about it. The difference in the two roles is the difference in the two intuitions.
\end{reply}

\begin{objection}[Metalinguistic negotiation]
The datum of Principle~\ref{p:datum} can be explained without sameness of meaning. The communities may be engaged in a metalinguistic negotiation about how \enquote{wrong} should be used, which is a genuine disagreement about a normative question even if their words differ in meaning \citep{PlunkettSundell2013}.
\end{objection}

\begin{reply}
Metalinguistic negotiation occurs, and the role semantics can allow that it occurs in the etiquette case, where the parties might argue about which of their words to adopt. But the explanation does not fit Horgan and Timmons's case. The disputants there take themselves to disagree about the act, not about the word, and they settle their dispute, when they do, by considering features of the act and its effects. A semantics that makes their dispute a dispute about words attributes to them a subject matter they do not recognize. The role semantics takes their dispute at face value, and Theorem~\ref{t:same} shows that it can.
\end{reply}

\begin{objection}[Circularity of the role]
The role semantics identifies the role of \enquote{wrong} by such features as being subject to correction by moral inquiry. That presupposes that we can identify moral inquiry independently of the meaning of moral terms.
\end{objection}

\begin{reply}
The features of the role are features of practice that can be identified without moral vocabulary: whether a community's use of a word is constrained by facts about acts, whether it guides action, whether it is subject to revision in the light of argument and evidence, whether its users treat it as applying to everyone alike, and whether they regard themselves as fallible about it. An anthropologist can identify these features in a community whose language she is learning, before she knows what its words mean. That is how she decides which of its words to translate as \enquote{wrong}, and the decision is the one that the role semantics makes for the Twin Earthers.
\end{reply}

\section{What the result does not establish}

The result does not establish which acts are wrong, and so it does not say which community, in Horgan and Timmons's case, is in error, or whether both are. It does not establish that every evaluative term of every community refers to a property, since the etiquette planet's term may refer to nothing more than conformity to custom. And it does not by itself establish the realization claim, that idealized disapproval is the unique realizer of the role; that is argued in MTE-01.

\section{Conclusion}

Moral Twin Earth refutes the causal-regulation semantics, as Horgan and Timmons argued. It does not refute naturalism. On the role semantics of MTE-01, moral terms refer to the property that realizes the role moral inquiry assigns them, and the Twin Earthers' inquiry assigns their terms the same role as ours. The terms therefore refer to the same property, the convergence of idealized disapproval, and the communities that dispute an act's wrongness genuinely disagree. Where their applications diverge, at least one is in error.

\begin{thebibliography}{99}\small
\bibitem[Boyd(1988)]{Boyd1988} Boyd, R. (1988). How to be a moral realist. In G. Sayre-McCord (ed.), \emph{Essays on Moral Realism}, 181--228. Cornell University Press.
\bibitem[Copp(2000)]{Copp2000} Copp, D. (2000). Milk, honey, and the good life on moral twin earth. \emph{Synthese}, 124(1), 113--137.
\bibitem[Dowell(2016)]{Dowell2016} Dowell, J.~L. (2016). The metaethical insignificance of moral twin earth. \emph{Oxford Studies in Metaethics}, 11.
\bibitem[Horgan and Timmons(1991)]{HorganTimmons1991} Horgan, T. and Timmons, M. (1991). New wave moral realism meets moral twin earth. \emph{Journal of Philosophical Research}, 16, 447--465.
\bibitem[Horgan and Timmons(1992a)]{HorganTimmons1992a} Horgan, T. and Timmons, M. (1992a). Troubles on moral twin earth: Moral queerness revived. \emph{Synthese}, 92(2), 221--260.
\bibitem[Horgan and Timmons(1992b)]{HorganTimmons1992b} Horgan, T. and Timmons, M. (1992b). Troubles for new wave moral semantics: The \enquote{open question argument} revived. \emph{Philosophical Papers}, 21(3), 153--175.
\bibitem[Merli(2002)]{Merli2002} Merli, D. (2002). Return to moral twin earth. \emph{Canadian Journal of Philosophy}, 32(2), 207--240.
\bibitem[Plunkett and Sundell(2013)]{PlunkettSundell2013} Plunkett, D. and Sundell, T. (2013). Disagreement and the semantics of normative and evaluative terms. \emph{Philosophers' Imprint}, 13(23).
\bibitem[Putnam(1975)]{Putnam1975} Putnam, H. (1975). The meaning of \enquote{meaning}. \emph{Minnesota Studies in the Philosophy of Science}, 7, 131--193.
\bibitem[van Roojen(2006)]{vanRoojen2006} van Roojen, M. (2006). Knowing enough to disagree: A new response to the moral twin earth argument. \emph{Oxford Studies in Metaethics}, 1.
\end{thebibliography}

\end{document}
